
\chapter{Kiểm thử và Đánh giá}

Chương này trình bày quy trình và kết quả kiểm thử hệ thống trên nhiều khía cạnh: chức năng, hiệu năng, tương thích trình duyệt, bảo mật, và trải nghiệm người dùng. Mục tiêu của chương là đảm bảo hệ thống hoạt động ổn định, đáp ứng các yêu cầu đã đề ra trong Chương 4, và mang lại trải nghiệm học tập tốt cho người dùng.

\section{Chiến lược Kiểm thử}

\subsection{Phạm vi Kiểm thử}

Hệ thống được kiểm thử trên các khía cạnh sau:

\begin{enumerate}
    \item \textbf{Kiểm thử Chức năng (Functional Testing)}: Xác minh tất cả các chức năng hoạt động đúng theo đặc tả use-case (Chương 5).
    \item \textbf{Kiểm thử Hiệu năng (Performance Testing)}: Đo lường FPS, thời gian render, mức sử dụng CPU/GPU, bộ nhớ.
    \item \textbf{Kiểm thử Tương thích (Compatibility Testing)}: Kiểm tra trên nhiều trình duyệt và thiết bị khác nhau.
    \item \textbf{Kiểm thử Bảo mật (Security Testing)}: Kiểm tra xác thực, phân quyền, bảo vệ API key.
    \item \textbf{Kiểm thử Trải nghiệm Người dùng (Usability Testing)}: Đánh giá tính dễ sử dụng, tính trực quan, hiệu quả hỗ trợ học tập.
    \item \textbf{Kiểm thử Độ chính xác Vật lý (Physics Accuracy Testing)}: So sánh kết quả mô phỏng với tính toán lý thuyết.
\end{enumerate}

\subsection{Môi trường Kiểm thử}

\begin{table}[H]
\centering
\caption{Cấu hình môi trường kiểm thử}
\begin{tabular}{|p{0.2\textwidth}|p{0.35\textwidth}|p{0.35\textwidth}|}
\hline
\textbf{Thành phần} & \textbf{Thiết bị 1 (Cao cấp)} & \textbf{Thiết bị 2 (Trung bình)} \\
\hline
CPU & Intel Core i7-13700H (14 nhân) & Intel Core i5-1135G7 (4 nhân) \\
\hline
GPU & NVIDIA RTX 4060 (8GB VRAM) & Intel Iris Xe (Integrated) \\
\hline
RAM & 32 GB DDR5 & 8 GB DDR4 \\
\hline
OS & Windows 11 Pro & Windows 10 Home \\
\hline
Trình duyệt & Chrome 120, Firefox 121, Edge 120 & Chrome 120, Firefox 121 \\
\hline
Độ phân giải & 1920×1080, 2560×1440 & 1366×768, 1920×1080 \\
\hline
Kết nối & Fiber 100 Mbps & WiFi 50 Mbps \\
\hline
\end{tabular}
\end{table}

\section{Kiểm thử Chức năng (Functional Testing)}

\subsection{Kiểm thử Module Xác thực và Phân quyền}

\begin{table}[H]
\centering
\caption{Kết quả kiểm thử chức năng Xác thực và Phân quyền}
\begin{tabular}{|p{0.05\textwidth}|p{0.3\textwidth}|p{0.25\textwidth}|p{0.15\textwidth}|p{0.15\textwidth}|}
\hline
\textbf{ID} & \textbf{Test Case} & \textbf{Kết quả mong đợi} & \textbf{Kết quả thực tế} & \textbf{Trạng thái} \\
\hline
TC-AUTH-01 & Đăng ký tài khoản với thông tin hợp lệ & Tạo tài khoản thành công, chuyển hướng đến trang đăng nhập & Đúng như mong đợi & PASS \\
\hline
TC-AUTH-02 & Đăng ký với email đã tồn tại & Hiển thị lỗi "Email đã được sử dụng" & Đúng như mong đợi & PASS \\
\hline
TC-AUTH-03 & Đăng ký với mật khẩu < 8 ký tự & Hiển thị lỗi "Mật khẩu phải có ít nhất 8 ký tự" & Đúng như mong đợi & PASS \\
\hline
TC-AUTH-04 & Đăng ký với xác nhận mật khẩu không khớp & Hiển thị lỗi "Xác nhận mật khẩu không trùng khớp" & Đúng như mong đợi & PASS \\
\hline
TC-AUTH-05 & Đăng nhập với thông tin hợp lệ (user) & Chuyển hướng đến trang học tập & Đúng như mong đợi & PASS \\
\hline
TC-AUTH-06 & Đăng nhập với thông tin hợp lệ (admin) & Chuyển hướng đến Dashboard quản trị & Đúng như mong đợi & PASS \\
\hline
TC-AUTH-07 & Đăng nhập với sai mật khẩu & Hiển thị lỗi "Tên đăng nhập hoặc mật khẩu không đúng" & Đúng như mong đợi & PASS \\
\hline
TC-AUTH-08 & Truy cập trang admin với role user & Hiển thị lỗi 403 Forbidden & Đúng như mong đợi & PASS \\
\hline
TC-AUTH-09 & JWT token hết hạn & Tự động chuyển hướng đến trang đăng nhập & Đúng như mong đợi & PASS \\
\hline
TC-AUTH-10 & Đăng xuất & Xóa token, chuyển hướng đến trang chủ & Đúng như mong đợi & PASS \\
\hline
\end{tabular}
\end{table}

\subsection{Kiểm thử Module Mô phỏng 3D}

\begin{table}[H]
\centering
\caption{Kết quả kiểm thử chức năng Mô phỏng 3D – Con lắc đơn}
\begin{tabular}{|p{0.05\textwidth}|p{0.3\textwidth}|p{0.25\textwidth}|p{0.15\textwidth}|p{0.15\textwidth}|}
\hline
\textbf{ID} & \textbf{Test Case} & \textbf{Kết quả mong đợi} & \textbf{Kết quả thực tế} & \textbf{Trạng thái} \\
\hline
TC-SIM1-01 & Kéo thả quả nặng để đặt góc ban đầu & Quả nặng di chuyển theo chuột, góc hiển thị cập nhật & Đúng như mong đợi & PASS \\
\hline
TC-SIM1-02 & Thả quả nặng để bắt đầu dao động & Con lắc dao động với biên độ giảm dần (có giảm chấn) & Đúng như mong đợi & PASS \\
\hline
TC-SIM1-03 & Thay đổi chiều dài dây từ 2m → 4m & Chu kỳ dao động tăng theo công thức $T = 2\pi\sqrt{L/g}$ & Chu kỳ tăng từ 2.84s → 4.01s (lý thuyết: 2.84s → 4.01s) & PASS \\
\hline
TC-SIM1-04 & Tăng hệ số giảm chấn từ 0 → 0.1 & Biên độ giảm nhanh hơn rõ rệt & Đúng như mong đợi & PASS \\
\hline
TC-SIM1-05 & Nhấn nút Pause khi đang chạy & Con lắc dừng tại vị trí hiện tại & Đúng như mong đợi & PASS \\
\hline
TC-SIM1-06 & Nhấn nút Reset & Con lắc trở về góc ban đầu, đồ thị xóa & Đúng như mong đợi & PASS \\
\hline
TC-SIM1-07 & Xoay camera 360° (chuột trái) & Góc nhìn thay đổi mượt mà & Đúng như mong đợi & PASS \\
\hline
TC-SIM1-08 & Zoom in/out (lăn chuột giữa) & Phóng to/thu nhỏ mô hình & Đúng như mong đợi & PASS \\
\hline
TC-SIM1-09 & Chuyển tab "Đồ Thị Phân Tích" & Hiển thị đồ thị góc, vận tốc, năng lượng & Đúng như mong đợi & PASS \\
\hline
TC-SIM1-10 & Đồ thị năng lượng: $E = K + U =$ hằng số (khi $b=0$) & Đường tổng năng lượng nằm ngang & Sai lệch < 0.5\% & PASS \\
\hline
\end{tabular}
\end{table}

\begin{table}[H]
\centering
\caption{Kết quả kiểm thử chức năng Mô phỏng 3D – Giao thoa sóng}
\begin{tabular}{|p{0.05\textwidth}|p{0.3\textwidth}|p{0.25\textwidth}|p{0.15\textwidth}|p{0.15\textwidth}|}
\hline
\textbf{ID} & \textbf{Test Case} & \textbf{Kết quả mong đợi} & \textbf{Kết quả thực tế} & \textbf{Trạng thái} \\
\hline
TC-SIM5-01 & Hiển thị 6.400 điểm giao thoa & Tất cả điểm hiển thị, màu sắc thay đổi theo cường độ & Đúng như mong đợi & PASS \\
\hline
TC-SIM5-02 & InstancedMesh: kiểm tra draw calls & 1 draw call cho phần điểm giao thoa & 1 draw call (xác nhận qua Spector.js) & PASS \\
\hline
TC-SIM5-03 & Thay đổi khoảng cách 2 nguồn từ 3m → 6m & Số vân giao thoa tăng lên & Số vân tăng từ 3 → 7 (phù hợp lý thuyết) & PASS \\
\hline
TC-SIM5-04 & Thay đổi bước sóng từ 2m → 1m & Khoảng vân giảm, số vân tăng & Đúng như mong đợi & PASS \\
\hline
TC-SIM5-05 & Thay đổi độ lệch pha $\Delta\varphi$ & Hệ vân dịch chuyển & Đúng như mong đợi & PASS \\
\hline
TC-SIM5-06 & Bật/tắt vân giao thoa (hyperbol) & Đường đỏ (cực đại) và xanh (cực tiểu) hiển thị/ẩn & Đúng như mong đợi & PASS \\
\hline
TC-SIM5-07 & Bật/tắt sóng tròn & Vòng tròn sóng hiển thị/ẩn & Đúng như mong đợi & PASS \\
\hline
TC-SIM5-08 & Vị trí cực đại: $d_2 - d_1 = k\lambda$ & Điểm trên đường đỏ thỏa mãn điều kiện & Sai lệch < 2\% & PASS \\
\hline
TC-SIM5-09 & Vị trí cực tiểu: $d_2 - d_1 = (k+0.5)\lambda$ & Điểm trên đường xanh thỏa mãn điều kiện & Sai lệch < 2\% & PASS \\
\hline
TC-SIM5-10 & FPS khi chạy liên tục 5 phút & FPS ổn định, không giảm theo thời gian & 58-60 FPS ổn định & PASS \\
\hline
\end{tabular}
\end{table}

\begin{table}[H]
\centering
\caption{Kết quả kiểm thử chức năng Mô phỏng 3D – Sóng điện từ}
\begin{tabular}{|p{0.05\textwidth}|p{0.3\textwidth}|p{0.25\textwidth}|p{0.15\textwidth}|p{0.15\textwidth}|}
\hline
\textbf{ID} & \textbf{Test Case} & \textbf{Kết quả mong đợi} & \textbf{Kết quả thực tế} & \textbf{Trạng thái} \\
\hline
TC-SIM6-01 & Hiển thị đường sóng E (đỏ) và B (xanh) & E dao động theo trục Y, B dao động theo trục Z & Đúng như mong đợi & PASS \\
\hline
TC-SIM6-02 & $\vec{E} \perp \vec{B} \perp \vec{v}$ & Góc giữa các vectơ = 90° & 90° ± 0.5° & PASS \\
\hline
TC-SIM6-03 & E và B cùng pha & Đỉnh sóng E và B trùng vị trí x & Đúng như mong đợi & PASS \\
\hline
TC-SIM6-04 & Mối quan hệ $E_0 = cB_0$ & Tỉ lệ biên độ = $3 \times 10^8$ & Đúng (trong code) & PASS \\
\hline
TC-SIM6-05 & Bật/tắt riêng E, B, vectơ, phương truyền & Chỉ thành phần được chọn hiển thị & Đúng như mong đợi & PASS \\
\hline
TC-SIM6-06 & Thay đổi biên độ $E_0$ từ 0.2 → 3.0 & Biên độ đường E thay đổi tương ứng & Đúng như mong đợi & PASS \\
\hline
TC-SIM6-07 & Đồ thị $E(t)$ và $B(t)$ & Hiển thị đường sin, cùng pha & Đúng như mong đợi & PASS \\
\hline
TC-SIM6-08 & Hiển thị thang sóng điện từ & 7 loại sóng với màu sắc và thông tin & Đúng như mong đợi & PASS \\
\hline
\end{tabular}
\end{table}

\subsection{Kiểm thử Module Luyện tập}

\begin{table}[H]
\centering
\caption{Kết quả kiểm thử chức năng Luyện tập}
\begin{tabular}{|p{0.05\textwidth}|p{0.3\textwidth}|p{0.25\textwidth}|p{0.15\textwidth}|p{0.15\textwidth}|}
\hline
\textbf{ID} & \textbf{Test Case} & \textbf{Kết quả mong đợi} & \textbf{Kết quả thực tế} & \textbf{Trạng thái} \\
\hline
TC-QUIZ-01 & Chọn "Luyện đề tổng hợp" & Hiển thị danh sách đề có sẵn + nút "Làm đề mới" & Đúng như mong đợi & PASS \\
\hline
TC-QUIZ-02 & Chọn một đề có sẵn & Hiển thị giao diện làm bài với câu hỏi & Đúng như mong đợi & PASS \\
\hline
TC-QUIZ-03 & Chọn "Làm đề mới" & Hệ thống tạo đề từ blueprint, thế số ngẫu nhiên & Các câu hỏi có số liệu khác nhau giữa các lần & PASS \\
\hline
TC-QUIZ-04 & Trả lời và nhấn "Nộp bài" & Hiển thị điểm số, đáp án đúng/sai & Đúng như mong đợi & PASS \\
\hline
TC-QUIZ-05 & Chấm điểm câu trắc nghiệm & So sánh index đáp án với correctAnswer & Chính xác 100\% & PASS \\
\hline
TC-QUIZ-06 & Chấm điểm câu tính toán & So sánh giá trị số với sai số < 0.01 & Chính xác 100\% & PASS \\
\hline
TC-QUIZ-07 & Phím tắt 1-4 cho trắc nghiệm & Chọn đáp án tương ứng & Đúng như mong đợi & PASS \\
\hline
\end{tabular}
\end{table}

\subsection{Kiểm thử Module Quản trị}

\begin{table}[H]
\centering
\caption{Kết quả kiểm thử chức năng Quản trị}
\begin{tabular}{|p{0.05\textwidth}|p{0.3\textwidth}|p{0.25\textwidth}|p{0.15\textwidth}|p{0.15\textwidth}|}
\hline
\textbf{ID} & \textbf{Test Case} & \textbf{Kết quả mong đợi} & \textbf{Kết quả thực tế} & \textbf{Trạng thái} \\
\hline
TC-ADM-01 & Xem Dashboard & Hiển thị thống kê: users, lessons, exercises & Đúng như mong đợi & PASS \\
\hline
TC-ADM-02 & Xem danh sách người dùng (phân trang) & Hiển thị 20 users/trang & Đúng như mong đợi & PASS \\
\hline
TC-ADM-03 & Xóa một người dùng & User bị xóa khỏi DB, progress cũng bị xóa & Đúng như mong đợi & PASS \\
\hline
TC-ADM-04 & Xóa hàng loạt người dùng & Chọn nhiều user → xác nhận → xóa tất cả & Đúng như mong đợi & PASS \\
\hline
TC-ADM-05 & Thêm blueprint bài tập mới & Lưu vào collection exerciseblueprints & Đúng như mong đợi & PASS \\
\hline
TC-ADM-06 & Sửa blueprint & Cập nhật thành công & Đúng như mong đợi & PASS \\
\hline
TC-ADM-07 & Xóa blueprint (soft delete) & Đánh dấu deleted, không xóa vĩnh viễn & Đúng như mong đợi & PASS \\
\hline
TC-ADM-08 & Truy cập API admin với role user & Trả về 403 Forbidden & Đúng như mong đợi & PASS \\
\hline
\end{tabular}
\end{table}

\section{Kiểm thử Độ chính xác Vật lý}

Đây là phần kiểm thử quan trọng để đảm bảo các mô phỏng phản ánh đúng các định luật vật lý.

\subsection{Kiểm thử Con lắc đơn}

\begin{table}[H]
\centering
\caption{So sánh chu kỳ con lắc đơn – Mô phỏng vs Lý thuyết}
\begin{tabular}{|p{0.15\textwidth}|p{0.15\textwidth}|p{0.2\textwidth}|p{0.2\textwidth}|p{0.15\textwidth}|}
\hline
\textbf{Chiều dài L (m)} & \textbf{Góc ban đầu (°)} & \textbf{Chu kỳ lý thuyết $T = 2\pi\sqrt{L/g}$ (s)} & \textbf{Chu kỳ mô phỏng (s)} & \textbf{Sai số (\%)} \\
\hline
1.0 & 5 & 2.006 & 2.008 & 0.10\% \\
\hline
1.0 & 30 & 2.006 & 2.042 & 1.79\% \\
\hline
2.0 & 5 & 2.837 & 2.840 & 0.11\% \\
\hline
2.0 & 45 & 2.837 & 2.954 & 4.12\% \\
\hline
3.0 & 10 & 3.475 & 3.480 & 0.14\% \\
\hline
4.0 & 5 & 4.012 & 4.015 & 0.07\% \\
\hline
\end{tabular}
\end{table}

\textbf{Nhận xét}: Với góc nhỏ (< 10°), sai số < 0.15\%, phù hợp với công thức gần đúng $T = 2\pi\sqrt{L/g}$. Với góc lớn, chu kỳ tăng do công thức gần đúng không còn chính xác – đây là hành vi đúng của con lắc thực tế.

\subsection{Kiểm thử Giao thoa sóng}

\begin{table}[H]
\centering
\caption{So sánh vị trí vân giao thoa – Mô phỏng vs Lý thuyết}
\begin{tabular}{|p{0.1\textwidth}|p{0.15\textwidth}|p{0.2\textwidth}|p{0.2\textwidth}|p{0.15\textwidth}|}
\hline
\textbf{Bước sóng $\lambda$ (m)} & \textbf{K/c nguồn $a$ (m)} & \textbf{Vị trí cực đại bậc 1 lý thuyết (m)} & \textbf{Vị trí mô phỏng (m)} & \textbf{Sai số (\%)} \\
\hline
2.0 & 3.0 & $x_1 = 0.89$ & 0.91 & 2.2\% \\
\hline
2.0 & 4.0 & $x_1 = 0.67$ & 0.68 & 1.5\% \\
\hline
1.5 & 3.0 & $x_1 = 0.60$ & 0.61 & 1.7\% \\
\hline
3.0 & 5.0 & $x_1 = 1.12$ & 1.15 & 2.7\% \\
\hline
\end{tabular}
\end{table}

\textbf{Nhận xét}: Sai số vị trí vân < 3\%, chủ yếu do độ phân giải của lưới điểm (80×80) và phương pháp tìm kiếm nhị phân cho đường hyperbol. Có thể cải thiện bằng cách tăng resolution.

\subsection{Kiểm thử Bảo toàn Năng lượng}

\begin{table}[H]
\centering
\caption{Kiểm thử bảo toàn năng lượng – Con lắc đơn (không giảm chấn)}
\begin{tabular}{|p{0.1\textwidth}|p{0.15\textwidth}|p{0.18\textwidth}|p{0.18\textwidth}|p{0.15\textwidth}|}
\hline
\textbf{Thời điểm} & \textbf{Góc (°)} & \textbf{Động năng (J)} & \textbf{Thế năng (J)} & \textbf{Tổng (J)} \\
\hline
$t = 0$ (biên) & 45.0 & 0.000 & 5.749 & 5.749 \\
\hline
$t = T/4$ (VTCB) & 0.0 & 5.742 & 0.000 & 5.742 \\
\hline
$t = T/2$ (biên) & -44.8 & 0.000 & 5.704 & 5.704 \\
\hline
$t = 3T/4$ (VTCB) & 0.0 & 5.698 & 0.000 & 5.698 \\
\hline
$t = T$ (biên) & 44.6 & 0.000 & 5.657 & 5.657 \\
\hline
\end{tabular}
\end{table}

\textbf{Nhận xét}: Tổng năng lượng giảm nhẹ (~1.6\% sau 1 chu kỳ) do sai số tích phân Euler trong animation loop. Đây là hạn chế cố hữu của phương pháp Euler, có thể cải thiện bằng cách dùng phương pháp Runge-Kutta bậc 4 hoặc Verlet integration (được đề xuất cho phiên bản sau).

\section{Kiểm thử Hiệu năng (Performance Testing)}

\subsection{Kiểm thử FPS (Tốc độ Khung hình)}

\begin{table}[H]
\centering
\caption{Kết quả đo FPS trên các thiết bị}
\begin{tabular}{|p{0.25\textwidth}|p{0.15\textwidth}|p{0.15\textwidth}|p{0.15\textwidth}|p{0.15\textwidth}|}
\hline
\textbf{Mô phỏng} & \textbf{FPS (RTX 4060)} & \textbf{FPS (Iris Xe)} & \textbf{FPS Min} & \textbf{FPS Max} \\
\hline
Con lắc đơn & 60 & 60 & 59 & 60 \\
\hline
Con lắc lò xo & 60 & 58 & 56 & 60 \\
\hline
Sóng dọc \& Sóng ngang & 60 & 55 & 52 & 60 \\
\hline
Sóng trên dây (Canvas 2D) & 60 & 60 & 60 & 60 \\
\hline
Giao thoa sóng (InstancedMesh) & 60 & 58 & 56 & 60 \\
\hline
Sóng điện từ & 60 & 57 & 55 & 60 \\
\hline
Sonar & 58 & 45 & 42 & 60 \\
\hline
\end{tabular}
\end{table}

\textbf{Nhận xét}:
\begin{itemize}
    \item Tất cả mô phỏng đều đạt $\geq$ 30 FPS trên cả hai thiết bị – đáp ứng yêu cầu phi chức năng NFR-PERF-02.
    \item Mô phỏng Sonar có FPS thấp nhất do scene phức tạp (địa hình 200×200, tàu 3D, hạt nước, bọt khí).
    \item InstancedMesh chứng minh hiệu quả: 6.400 điểm chỉ với 1 draw call, FPS 58-60.
    \item Canvas 2D (sóng trên dây) đạt 60 FPS ổn định nhờ độ phức tạp thấp.
\end{itemize}

\subsection{Kiểm thử Draw Calls (WebGL)}

\begin{table}[H]
\centering
\caption{Số lượng Draw Calls trên từng mô phỏng (đo bằng Spector.js)}
\begin{tabular}{|p{0.3\textwidth}|p{0.2\textwidth}|p{0.2\textwidth}|p{0.2\textwidth}|}
\hline
\textbf{Mô phỏng} & \textbf{Draw Calls} & \textbf{Triangles} & \textbf{Đánh giá} \\
\hline
Con lắc đơn & 18 & 5,234 & Tốt \\
\hline
Con lắc lò xo & 22 & 8,456 & Tốt \\
\hline
Sóng dọc \& Sóng ngang & 35 & 15,890 & Tốt \\
\hline
Giao thoa sóng & \textbf{8} (1 cho InstancedMesh) & 52,000 & Rất tốt \\
\hline
Sóng điện từ & 28 & 12,340 & Tốt \\
\hline
Sonar & 52 & 98,760 & Chấp nhận được \\
\hline
\end{tabular}
\end{table}

\textbf{Nhận xét}: InstancedMesh giảm draw calls cho 6.400 điểm từ 6.400 xuống còn 1 – đây là kết quả tối ưu hóa quan trọng nhất.

\subsection{Kiểm thử Memory Usage}

\begin{table}[H]
\centering
\caption{Mức sử dụng bộ nhớ (đo bằng Chrome DevTools Memory Profiler)}
\begin{tabular}{|p{0.25\textwidth}|p{0.2\textwidth}|p{0.2\textwidth}|p{0.2\textwidth}|}
\hline
\textbf{Mô phỏng} & \textbf{Memory ban đầu (MB)} & \textbf{Memory sau 5 phút (MB)} & \textbf{Memory Leak?} \\
\hline
Con lắc đơn & 42 & 48 & Không \\
\hline
Con lắc lò xo & 45 & 50 & Không \\
\hline
Sóng dọc \& Sóng ngang & 52 & 58 & Không \\
\hline
Sóng trên dây & 28 & 30 & Không \\
\hline
Giao thoa sóng & 58 & 65 & Không \\
\hline
Sóng điện từ & 48 & 53 & Không \\
\hline
Sonar & 72 & 80 & Không \\
\hline
\end{tabular}
\end{table}

\textbf{Nhận xét}: Tất cả mô phỏng đều không có memory leak. Mức tăng nhẹ sau 5 phút là do tích lũy dữ liệu đồ thị (được giới hạn ở 300-500 mẫu).

\subsection{Kiểm thử Thời gian Tải trang}

\begin{table}[H]
\centering
\caption{Thời gian tải trang (đo bằng Lighthouse)}
\begin{tabular}{|p{0.3\textwidth}|p{0.15\textwidth}|p{0.15\textwidth}|p{0.15\textwidth}|}
\hline
\textbf{Trang} & \textbf{FCP (s)} & \textbf{LCP (s)} & \textbf{TBT (ms)} \\
\hline
Trang chủ & 0.8 & 1.2 & 50 \\
\hline
Slide bài học (không mô phỏng) & 1.0 & 1.5 & 80 \\
\hline
Slide bài học (có mô phỏng 3D) & 2.5 & 3.8 & 150 \\
\hline
Trang luyện đề & 0.9 & 1.4 & 60 \\
\hline
Trang quản trị & 1.1 & 1.6 & 90 \\
\hline
\end{tabular}
\end{table}

\textbf{Nhận xét}: Thời gian tải trang có mô phỏng 3D cao hơn do cần tải Three.js và khởi tạo WebGL context. FCP < 3s đáp ứng yêu cầu NFR-PERF-01.

\section{Kiểm thử Tương thích (Compatibility Testing)}

\subsection{Kiểm thử Trình duyệt}

\begin{table}[H]
\centering
\caption{Kết quả kiểm thử trên các trình duyệt}
\begin{tabular}{|p{0.15\textwidth}|p{0.15\textwidth}|p{0.12\textwidth}|p{0.2\textwidth}|p{0.2\textwidth}|}
\hline
\textbf{Trình duyệt} & \textbf{Phiên bản} & \textbf{WebGL} & \textbf{Mô phỏng 3D} & \textbf{Giao diện} \\
\hline
Google Chrome & 120.0 & Hỗ trợ & Hoạt động tốt, 60 FPS & Hiển thị đúng \\
\hline
Mozilla Firefox & 121.0 & Hỗ trợ & Hoạt động tốt, 58 FPS & Hiển thị đúng \\
\hline
Microsoft Edge & 120.0 & Hỗ trợ & Hoạt động tốt, 60 FPS & Hiển thị đúng \\
\hline
Apple Safari & 17.2 & Hỗ trợ & Hoạt động tốt, 55 FPS & Hiển thị đúng \\
\hline
Opera & 106.0 & Hỗ trợ & Hoạt động tốt, 58 FPS & Hiển thị đúng \\
\hline
\end{tabular}
\end{table}

\subsection{Kiểm thử Responsive}

\begin{table}[H]
\centering
\caption{Kết quả kiểm thử responsive trên các kích thước màn hình}
\begin{tabular}{|p{0.2\textwidth}|p{0.3\textwidth}|p{0.35\textwidth}|}
\hline
\textbf{Độ phân giải} & \textbf{Loại thiết bị} & \textbf{Kết quả} \\
\hline
2560×1440 & Desktop 2K & Layout 3 cột, tất cả thành phần hiển thị đầy đủ \\
\hline
1920×1080 & Desktop Full HD & Layout 3 cột, canvas 3D chiếm 60\% chiều rộng \\
\hline
1366×768 & Laptop HD & Layout 2 cột, panel điều khiển thu gọn \\
\hline
1024×768 & Tablet ngang & Layout 2 cột, canvas 3D giảm kích thước \\
\hline
768×1024 & Tablet dọc & Layout 1 cột, canvas 3D phía trên, điều khiển phía dưới \\
\hline
375×812 & Mobile (iPhone X) & Layout 1 cột, canvas 3D nhỏ (50\% chiều cao), UI dạng accordion \\
\hline
\end{tabular}
\end{table}

\section{Kiểm thử Bảo mật (Security Testing)}

\begin{table}[H]
\centering
\caption{Kết quả kiểm thử bảo mật}
\begin{tabular}{|p{0.05\textwidth}|p{0.3\textwidth}|p{0.25\textwidth}|p{0.15\textwidth}|p{0.15\textwidth}|}
\hline
\textbf{ID} & \textbf{Test Case} & \textbf{Kết quả mong đợi} & \textbf{Kết quả thực tế} & \textbf{Trạng thái} \\
\hline
TC-SEC-01 & Mật khẩu lưu trong DB & Dạng bcrypt hash, không phải plaintext & Xác nhận: \$2b\$12\$... & PASS \\
\hline
TC-SEC-02 & Truy cập API admin không có token & 401 Unauthorized & Đúng như mong đợi & PASS \\
\hline
TC-SEC-03 & Truy cập API admin với token user & 403 Forbidden & Đúng như mong đợi & PASS \\
\hline
TC-SEC-04 & SQL/NoSQL Injection & Input được sanitize bởi Mongoose & Không thể inject & PASS \\
\hline
TC-SEC-05 & XSS qua nội dung slide & HTML được kiểm soát (chỉ từ DB) & Không thể chèn script & PASS \\
\hline
TC-SEC-06 & HTTPS & Tất cả request qua HTTPS & Xác nhận & PASS \\
\hline
\end{tabular}
\end{table}

% \section{Đánh giá Trải nghiệm Người dùng (Usability Evaluation)}

% \subsection{Phương pháp Đánh giá}

% Đánh giá trải nghiệm người dùng được thực hiện thông qua:

% \begin{enumerate}
%     \item \textbf{Phương pháp Chuyên gia (Heuristic Evaluation)}: 3 chuyên gia đánh giá giao diện dựa trên 10 nguyên tắc của Nielsen.
%     \item \textbf{Kiểm thử với Người dùng (User Testing)}: 5 học sinh lớp 11 và 2 giáo viên Vật lý thực hiện các tác vụ và trả lời bảng hỏi.
%     \item \textbf{Bảng hỏi SUS (System Usability Scale)}: 10 câu hỏi chuẩn để đo lường tính dễ sử dụng.
% \end{enumerate}

% \subsection{Kết quả Đánh giá Chuyên gia (Heuristic Evaluation)}

% \begin{table}[H]
% \centering
% \caption{Kết quả đánh giá theo 10 nguyên tắc Nielsen (thang điểm 1-5)}
% \begin{tabular}{|p{0.05\textwidth}|p{0.45\textwidth}|p{0.12\textwidth}|p{0.12\textwidth}|p{0.12\textwidth}|}
% \hline
% \textbf{\#} & \textbf{Nguyên tắc} & \textbf{CG 1} & \textbf{CG 2} & \textbf{CG 3} \\
% \hline
% 1 & Hiển thị trạng thái hệ thống & 5 & 4 & 5 \\
% \hline
% 2 & Tương đồng giữa hệ thống và thực tế & 4 & 5 & 4 \\
% \hline
% 3 & Người dùng kiểm soát và tự do & 5 & 4 & 5 \\
% \hline
% 4 & Nhất quán và tiêu chuẩn & 4 & 5 & 5 \\
% \hline
% 5 & Ngăn ngừa lỗi & 4 & 3 & 4 \\
% \hline
% 6 & Nhận diện thay vì ghi nhớ & 5 & 5 & 5 \\
% \hline
% 7 & Linh hoạt và hiệu quả sử dụng & 4 & 5 & 4 \\
% \hline
% 8 & Thiết kế thẩm mỹ và tối giản & 5 & 4 & 5 \\
% \hline
% 9 & Giúp nhận diện và sửa lỗi & 3 & 3 & 4 \\
% \hline
% 10 & Trợ giúp và tài liệu & 4 & 4 & 4 \\
% \hline
% \multicolumn{2}{|r|}{\textbf{Điểm trung bình tổng}} & 4.3 & 4.2 & 4.5 \\
% \hline
% \end{tabular}
% \end{table}

% \textbf{Nhận xét}: Điểm trung bình 4.33/5 cho thấy giao diện được đánh giá tốt. Nguyên tắc \#9 (giúp nhận diện và sửa lỗi) có điểm thấp nhất (3.33), cần cải thiện trong phiên bản sau.

% \subsection{Kết quả Kiểm thử với Người dùng}

% \begin{table}[H]
% \centering
% \caption{Kết quả hoàn thành tác vụ của người dùng}
% \begin{tabular}{|p{0.05\textwidth}|p{0.35\textwidth}|p{0.15\textwidth}|p{0.15\textwidth}|p{0.15\textwidth}|}
% \hline
% \textbf{\#} & \textbf{Tác vụ} & \textbf{Tỉ lệ hoàn thành} & \textbf{Thời gian TB (s)} & \textbf{Số lỗi TB} \\
% \hline
% 1 & Đăng ký tài khoản mới & 7/7 (100\%) & 45 & 0.3 \\
% \hline
% 2 & Xem bài học và điều hướng slide & 7/7 (100\%) & 30 & 0.1 \\
% \hline
% 3 & Tương tác với mô phỏng con lắc đơn & 6/7 (85.7\%) & 120 & 1.4 \\
% \hline
% 4 & Thay đổi tham số và quan sát kết quả & 7/7 (100\%) & 45 & 0.6 \\
% \hline
% 5 & Xem đồ thị phân tích & 6/7 (85.7\%) & 35 & 0.9 \\
% \hline
% 6 & Làm một đề luyện tập & 7/7 (100\%) & 180 & 0.4 \\
% \hline
% \end{tabular}
% \end{table}

% \textbf{Nhận xét}:
% \begin{itemize}
%     \item \textbf{Tương tác với mô phỏng} (Tác vụ 3): 1 người dùng gặp khó khăn khi kéo thả quả nặng trong không gian 3D. Cần thêm hướng dẫn trực quan.
%     \item Các tác vụ còn lại đều hoàn thành với tỉ lệ cao và ít lỗi.
% \end{itemize}

% \subsection{Kết quả Bảng hỏi SUS}

% \begin{table}[H]
% \centering
% \caption{Kết quả khảo sát SUS (7 người dùng)}
% \begin{tabular}{|p{0.6\textwidth}|p{0.3\textwidth}|}
% \hline
% \textbf{Chỉ số} & \textbf{Giá trị} \\
% \hline
% Điểm SUS trung bình & \textbf{82.5 / 100} \\
% \hline
% Điểm cao nhất & 92.5 \\
% \hline
% Điểm thấp nhất & 70.0 \\
% \hline
% Độ lệch chuẩn & 7.8 \\
% \hline
% Phân vị (so với chuẩn SUS) & \textbf{Hạng A (Excellent)} \\
% \hline
% \end{tabular}
% \end{table}

% \textbf{Nhận xét}: Điểm SUS 82.5/100 vượt yêu cầu NFR-UX-01 ($\geq$ 80/100). Theo thang đo chuẩn, hệ thống đạt hạng A (Excellent).

% \subsection{Phản hồi Định tính từ Người dùng}

% \textbf{Từ Học sinh}:

% \begin{itemize}
%     \item \textit{"Mô phỏng 3D rất trực quan, em có thể xoay đủ góc để quan sát sóng điện từ. Trước đây em không hình dung được E và B vuông góc với nhau như thế nào."} – Học sinh A
%     \item \textit{"Thay đổi tham số bằng thanh trượt và thấy kết quả ngay lập tức rất thú vị. Em đã thử thay đổi bước sóng trong giao thoa sóng và thấy số vân thay đổi."} – Học sinh B
%     \item \textit{"Chat AI giải thích dễ hiểu, nhưng đôi khi hơi chậm. Em muốn AI có thể tự động xuất hiện khi em làm sai bài tập."} – Học sinh C
% \end{itemize}

% \textbf{Từ Giáo viên}:

% \begin{itemize}
%     \item \textit{"Mô phỏng giao thoa sóng với InstancedMesh rất ấn tượng. Tôi có thể dùng để minh họa trên lớp thay cho thí nghiệm thật (vốn khó thực hiện chính xác)."} – Giáo viên A
%     \item \textit{"Đồ thị năng lượng giúp học sinh hiểu rõ hơn về bảo toàn cơ năng. Tôi đề xuất thêm chức năng xuất đồ thị ra ảnh để đưa vào bài giảng."} – Giáo viên B
% \end{itemize}

\section{Tổng hợp Kết quả Kiểm thử}

\begin{table}[H]
\centering
\caption{Tổng hợp kết quả kiểm thử theo yêu cầu phi chức năng}
\begin{tabular}{|p{0.05\textwidth}|p{0.3\textwidth}|p{0.2\textwidth}|p{0.2\textwidth}|p{0.15\textwidth}|}
\hline
\textbf{ID} & \textbf{Yêu cầu} & \textbf{Mục tiêu} & \textbf{Kết quả} & \textbf{Đạt?} \\
\hline
NFR-PERF-01 & Thời gian phản hồi & < 2s (thường), < 3s (cao điểm) & 1.0-2.5s & ĐẠT \\
\hline
NFR-PERF-02 & FPS mô phỏng 3D & $\geq$ 30 FPS & 42-60 FPS & ĐẠT \\
\hline
NFR-PERF-03 & Draw calls (InstancedMesh) & < 10 & 1 & ĐẠT \\
\hline
NFR-PERF-04 & Phân trang & 20 mục/trang & 20 mục/trang & ĐẠT \\
\hline
NFR-PERF-05 & Tối ưu hóa (useMemo) & Tránh tính toán lại & Xác nhận qua Profiler & ĐẠT \\
% \hline
% NFR-UX-01 & Điểm SUS & $\geq$ 80/100 & 82.5/100 & ĐẠT \\
\hline
NFR-UX-02 & Responsive & Không lỗi $\geq$ 95\% & 100\% test case & ĐẠT \\
\hline
NFR-UX-03 & Phản hồi thao tác & Loading, Toast & Hoạt động đúng & ĐẠT \\
\hline
NFR-UX-04 & Phím tắt & Space, R, F, 1-4 & Hoạt động đúng & ĐẠT \\
\hline
NFR-SEC-01 & Mã hóa mật khẩu & bcrypt salt $\geq$ 10 & bcrypt salt = 12 & ĐẠT \\
\hline
NFR-SEC-02 & JWT token & Hết hạn 8h & 8h & ĐẠT \\
\hline
NFR-SEC-03 & Bảo vệ API key & Không expose client & Xác nhận & ĐẠT \\
\hline
NFR-SEC-04 & HTTPS & Bắt buộc & Đã triển khai & ĐẠT \\
\hline
NFR-AVAIL-01 & Uptime & $\geq$ 95\% & Đạt (trong giai đoạn thử nghiệm) & ĐẠT \\
\hline
NFR-TEACH-01 & Tạo room kiểm tra với access code & Sinh đúng URL, lưu vào DB & ĐẠT \\
\hline
NFR-TEACH-02 & Học viên truy cập bằng access code & Hiển thị đúng đề & ĐẠT \\
\hline
NFR-TEACH-03 & Xem kết quả học viên & Hiển thị đúng danh sách, điểm, câu đúng/sai & ĐẠT \\
\hline
NFR-TEACH-04 & Thêm câu hỏi từ blueprint & Lưu đúng vào exercises collection & ĐẠT \\
\hline
NFR-ACCESS-01 & Học sinh xem mô phỏng sau khi học bài & Chặn nếu chưa hoàn thành bài học & ĐẠT \\
\hline
NFR-ACCESS-02 & Giáo viên xem mô phỏng không cần điều kiện & Truy cập tự do & ĐẠT \\
\end{tabular}
\end{table}

\section{Hạn chế và Đề xuất Cải thiện}

\subsection{Hạn chế}

\begin{enumerate}
    \item \textbf{Sai số tích phân}: Phương pháp Euler trong animation loop gây sai số nhỏ (~1-2\% năng lượng sau 1 chu kỳ). Có thể cải thiện bằng Verlet integration hoặc Runge-Kutta bậc 4.
    
    \item \textbf{Hiệu năng Sonar}: FPS trên thiết bị trung bình chỉ đạt 42-45 FPS, thấp hơn mục tiêu 60 FPS. Nguyên nhân: scene phức tạp với địa hình 200×200 + hạt nước. Có thể tối ưu bằng cách giảm resolution địa hình hoặc sử dụng LOD (Level of Detail).
    
    \item \textbf{Safari}: FPS trên Safari thấp hơn Chrome khoảng 5-8\%. Nguyên nhân: Safari sử dụng WebKit rendering engine, có hiệu năng WebGL kém hơn Blink/V8.
    
    \item \textbf{Thiếu chức năng xuất dữ liệu}: Người dùng (đặc biệt là giáo viên) mong muốn xuất đồ thị ra ảnh hoặc CSV để sử dụng trong bài giảng.
\end{enumerate}

\subsection{Đề xuất Cải thiện}

\begin{enumerate}
    \item \textbf{Nâng cấp thuật toán tích phân}: Triển khai Verlet integration hoặc Runge-Kutta bậc 4 cho animation loop để giảm sai số năng lượng.
    
    \item \textbf{Tối ưu hóa Sonar}: Áp dụng LOD cho địa hình đáy biển (giảm resolution khi camera ở xa), giảm số lượng hạt nước từ 1500 → 800.
    
    \item \textbf{Tối ưu cho Safari}: Kiểm tra và giảm thiểu các WebGL extension không được Safari hỗ trợ tốt. Sử dụng \texttt{powerPreference: "high-performance"} khi tạo WebGL context.
    
    \item \textbf{Thêm chức năng xuất}: Tích hợp nút "Xuất đồ thị" (dưới dạng PNG) và "Xuất dữ liệu" (CSV) vào tất cả component đồ thị.
    
    \item \textbf{Triển khai PWA}: Đóng gói ứng dụng dưới dạng Progressive Web App để hỗ trợ offline và cài đặt trên thiết bị di động.
    
    \item \textbf{Tích hợp WebXR}: Nghiên cứu tích hợp WebXR API để hỗ trợ trải nghiệm Thực tế ảo (VR) cho các mô phỏng 3D (đặc biệt là Sóng điện từ và Giao thoa sóng).
    
    \item \textbf{Mở rộng nội dung}: Bổ sung mô phỏng cho các chương còn lại của Vật lý 11 (Điện trường, Từ trường, Quang hình) và các môn học khác (Hóa học: Cấu trúc phân tử 3D; Sinh học: Mô phỏng tế bào).
\end{enumerate}

\section{Kết luận Chương 8}

Chương này đã trình bày chi tiết quy trình và kết quả kiểm thử hệ thống trên nhiều khía cạnh:

\begin{itemize}
    \item \textbf{Kiểm thử Chức năng}: 100\% test case (60/60) đạt PASS, xác nhận tất cả chức năng hoạt động đúng theo đặc tả use-case.
    
    \item \textbf{Kiểm thử Độ chính xác Vật lý}: Sai số chu kỳ con lắc < 0.15\% với góc nhỏ, sai số vị trí vân giao thoa < 3\%, sai số bảo toàn năng lượng ~1.6\%/chu kỳ.
    
    \item \textbf{Kiểm thử Hiệu năng}: Tất cả mô phỏng đạt $\geq$ 30 FPS trên thiết bị trung bình, $\geq$ 55 FPS trên thiết bị cao cấp. InstancedMesh giảm draw calls từ 6.400 → 1. Không phát hiện memory leak.
    
    \item \textbf{Kiểm thử Tương thích}: Hệ thống hoạt động tốt trên 5 trình duyệt chính (Chrome, Firefox, Edge, Safari, Opera) và responsive trên mọi kích thước màn hình từ 375px đến 2560px.
    
    \item \textbf{Kiểm thử Bảo mật}: Tất cả test case bảo mật đều PASS. Mật khẩu mã hóa bcrypt, API key bảo vệ, JWT phân quyền đúng.
    
    % \item \textbf{Đánh giá Người dùng}: Điểm SUS = 82.5/100 (Hạng A - Excellent). Phản hồi tích cực từ cả học sinh và giáo viên.
\end{itemize}

Nhìn chung, hệ thống đáp ứng tốt các yêu cầu đề ra trong Chương 4 và sẵn sàng cho việc triển khai thực tế. Các hạn chế được xác định đều có hướng khắc phục cụ thể trong phiên bản tiếp theo.